% Options for packages loaded elsewhere
% Options for packages loaded elsewhere
\PassOptionsToPackage{unicode}{hyperref}
\PassOptionsToPackage{hyphens}{url}
\PassOptionsToPackage{dvipsnames,svgnames,x11names}{xcolor}
%
\documentclass[
  russian,
]{article}
\usepackage{xcolor}
\usepackage[top=20mm,left=20mm,right=20mm,bottom=20mm]{geometry}
\usepackage{amsmath,amssymb}
\setcounter{secnumdepth}{5}
\usepackage{iftex}
\ifPDFTeX
  \usepackage[T1]{fontenc}
  \usepackage[utf8]{inputenc}
  \usepackage{textcomp} % provide euro and other symbols
\else % if luatex or xetex
  \usepackage{unicode-math} % this also loads fontspec
  \defaultfontfeatures{Scale=MatchLowercase}
  \defaultfontfeatures[\rmfamily]{Ligatures=TeX,Scale=1}
\fi
\usepackage{lmodern}
\ifPDFTeX\else
  % xetex/luatex font selection
  \setmainfont[]{Times New Roman}
  \setsansfont[]{Arial}
  \setmonofont[]{Courier New}
\fi
% Use upquote if available, for straight quotes in verbatim environments
\IfFileExists{upquote.sty}{\usepackage{upquote}}{}
\IfFileExists{microtype.sty}{% use microtype if available
  \usepackage[]{microtype}
  \UseMicrotypeSet[protrusion]{basicmath} % disable protrusion for tt fonts
}{}
\makeatletter
\@ifundefined{KOMAClassName}{% if non-KOMA class
  \IfFileExists{parskip.sty}{%
    \usepackage{parskip}
  }{% else
    \setlength{\parindent}{0pt}
    \setlength{\parskip}{6pt plus 2pt minus 1pt}}
}{% if KOMA class
  \KOMAoptions{parskip=half}}
\makeatother
% Make \paragraph and \subparagraph free-standing
\makeatletter
\ifx\paragraph\undefined\else
  \let\oldparagraph\paragraph
  \renewcommand{\paragraph}{
    \@ifstar
      \xxxParagraphStar
      \xxxParagraphNoStar
  }
  \newcommand{\xxxParagraphStar}[1]{\oldparagraph*{#1}\mbox{}}
  \newcommand{\xxxParagraphNoStar}[1]{\oldparagraph{#1}\mbox{}}
\fi
\ifx\subparagraph\undefined\else
  \let\oldsubparagraph\subparagraph
  \renewcommand{\subparagraph}{
    \@ifstar
      \xxxSubParagraphStar
      \xxxSubParagraphNoStar
  }
  \newcommand{\xxxSubParagraphStar}[1]{\oldsubparagraph*{#1}\mbox{}}
  \newcommand{\xxxSubParagraphNoStar}[1]{\oldsubparagraph{#1}\mbox{}}
\fi
\makeatother


\usepackage{longtable,booktabs,array}
\newcounter{none} % for unnumbered tables
\usepackage{calc} % for calculating minipage widths
% Correct order of tables after \paragraph or \subparagraph
\usepackage{etoolbox}
\makeatletter
\patchcmd\longtable{\par}{\if@noskipsec\mbox{}\fi\par}{}{}
\makeatother
% Allow footnotes in longtable head/foot
\IfFileExists{footnotehyper.sty}{\usepackage{footnotehyper}}{\usepackage{footnote}}
\makesavenoteenv{longtable}
\usepackage{graphicx}
\makeatletter
\newsavebox\pandoc@box
\newcommand*\pandocbounded[1]{% scales image to fit in text height/width
  \sbox\pandoc@box{#1}%
  \Gscale@div\@tempa{\textheight}{\dimexpr\ht\pandoc@box+\dp\pandoc@box\relax}%
  \Gscale@div\@tempb{\linewidth}{\wd\pandoc@box}%
  \ifdim\@tempb\p@<\@tempa\p@\let\@tempa\@tempb\fi% select the smaller of both
  \ifdim\@tempa\p@<\p@\scalebox{\@tempa}{\usebox\pandoc@box}%
  \else\usebox{\pandoc@box}%
  \fi%
}
% Set default figure placement to htbp
\def\fps@figure{htbp}
\makeatother



\ifLuaTeX
\usepackage[bidi=basic,shorthands=off,provide=*]{babel}
\else
\usepackage[bidi=default,shorthands=off,provide=*]{babel}
\fi
\ifPDFTeX
\else
\babelfont{rm}[]{Times New Roman}
\fi
\ifLuaTeX
  \usepackage{selnolig} % disable illegal ligatures
\fi


\setlength{\emergencystretch}{3em} % prevent overfull lines

\providecommand{\tightlist}{%
  \setlength{\itemsep}{0pt}\setlength{\parskip}{0pt}}



 


\makeatletter
\@ifpackageloaded{caption}{}{\usepackage{caption}}
\AtBeginDocument{%
\ifdefined\contentsname
  \renewcommand*\contentsname{Содержание}
\else
  \newcommand\contentsname{Содержание}
\fi
\ifdefined\listfigurename
  \renewcommand*\listfigurename{Список Иллюстраций}
\else
  \newcommand\listfigurename{Список Иллюстраций}
\fi
\ifdefined\listtablename
  \renewcommand*\listtablename{Список Таблиц}
\else
  \newcommand\listtablename{Список Таблиц}
\fi
\ifdefined\figurename
  \renewcommand*\figurename{Рисунок}
\else
  \newcommand\figurename{Рисунок}
\fi
\ifdefined\tablename
  \renewcommand*\tablename{Таблица}
\else
  \newcommand\tablename{Таблица}
\fi
}
\@ifpackageloaded{float}{}{\usepackage{float}}
\floatstyle{ruled}
\@ifundefined{c@chapter}{\newfloat{codelisting}{h}{lop}}{\newfloat{codelisting}{h}{lop}[chapter]}
\floatname{codelisting}{Список}
\newcommand*\listoflistings{\listof{codelisting}{Список Каталогов}}
\makeatother
\makeatletter
\makeatother
\makeatletter
\@ifpackageloaded{caption}{}{\usepackage{caption}}
\@ifpackageloaded{subcaption}{}{\usepackage{subcaption}}
\makeatother
\usepackage{bookmark}
\IfFileExists{xurl.sty}{\usepackage{xurl}}{} % add URL line breaks if available
\urlstyle{same}
\hypersetup{
  pdftitle={Задачи по Эконометрике-2: Предельные значения},
  pdfauthor={Н.В. Артамонов, А. Н. Лата (МГИМО МИД России)},
  pdflang={ru},
  colorlinks=true,
  linkcolor={blue},
  filecolor={Maroon},
  citecolor={Blue},
  urlcolor={Blue},
  pdfcreator={LaTeX via pandoc}}


\title{Задачи по Эконометрике-2: Предельные значения}
\author{Н.В. Артамонов, А. Н. Лата (МГИМО МИД России)}
\date{}
\begin{document}
\maketitle

\renewcommand*\contentsname{Содержание}
{
\setcounter{tocdepth}{3}
\tableofcontents
}

\section{Предельные
значения}\label{ux43fux440ux435ux434ux435ux43bux44cux43dux44bux435-ux437ux43dux430ux447ux435ux43dux438ux44f}

\subsection{для probit}\label{ux434ux43bux44f-probit}

Рассмотрим probit-модель \(P(y=1)=\Phi(x'\beta)\), где
\(x'\beta=\beta_0+\beta_1x_1+\dots+\beta_kx_k\)

Предельные значения
\(\frac{\partial P(y=1)}{\partial x_j}=\phi(x'\beta)\beta_j\), где
\(\phi(z)=\frac{1}{\sqrt{2\pi}}\exp\{-z^2/2\}\)

\subsection{для logit}\label{ux434ux43bux44f-logit}

Рассмотрим logit-модель \(P(y=1)=\Lambda(x'\beta)\), где
\(x'\beta=\beta_0+\beta_1x_1+\dots+\beta_kx_k\)

Предельные значения
\(\frac{\partial P(y=1)}{\partial x_j}=\lambda(x'\beta)\beta_j\), где
\(\lambda(z)=\frac{\exp(z)}{(1+\exp(z))^2}\)

\subsection{Средние предельные
значения}\label{ux441ux440ux435ux434ux43dux438ux435-ux43fux440ux435ux434ux435ux43bux44cux43dux44bux435-ux437ux43dux430ux447ux435ux43dux438ux44f}

Обычно рассчитываются

\begin{itemize}
\tightlist
\item
  предельные значения для каждого регрессора в средней точке (MEM -
  Marginal Effects at the Mean):

  \begin{itemize}
  \tightlist
  \item
    \(\phi(\bar{x}'\beta)\beta_j\) для probit
  \item
    \(\lambda(\bar{x}'\beta)\beta_j\) для logit
  \end{itemize}
\item
  среднее по всей выборке предельное значения для каждого регрессора
  (AME - Average Marginal Effects):

  \begin{itemize}
  \tightlist
  \item
    \(\overline{\phi(x'\beta)\beta_j}\) для probit
  \item
    \(\overline{\lambda(x'\beta)\beta_j}\) для logit
  \end{itemize}
\end{itemize}

\section{labour force equation}\label{labour-force-equation}

Для датасета \texttt{TableF5-1} рассмотрим регрессию \textbf{LFP на WA,
np.log(FAMINC), WE, KL6, K618, CIT, UN} трёх спецификаций:

\begin{itemize}
\tightlist
\item
  LPM
\item
  logit
\item
  probit
\end{itemize}

Результаты подгонки:

{\def\LTcaptype{none} % do not increment counter
\begin{longtable}[]{@{}llll@{}}
\toprule\noalign{}
& LPM & logit & probit \\
\midrule\noalign{}
\endhead
\bottomrule\noalign{}
\endlastfoot
Dep. Variable & LFP & LFP & LFP \\
WA & -0.013*** & -0.063*** & -0.038*** \\
& (0.003) & (0.013) & (0.008) \\
np.log(FAMINC) & 0.075** & 0.341** & 0.205** \\
& (0.037) & (0.172) & (0.104) \\
WE & 0.038*** & 0.179*** & 0.108*** \\
& (0.008) & (0.040) & (0.024) \\
KL6 & -0.302*** & -1.443*** & -0.868*** \\
& (0.036) & (0.194) & (0.112) \\
K618 & -0.018 & -0.095 & -0.057 \\
& (0.014) & (0.067) & (0.040) \\
CIT & -0.048 & -0.214 & -0.126 \\
& (0.038) & (0.176) & (0.107) \\
UN & -0.004 & -0.017 & -0.011 \\
& (0.006) & (0.026) & (0.016) \\
Intercept & 0.079 & -1.856 & -1.108 \\
& (0.356) & (1.679) & (1.014) \\
R-squared & 0.130 & & \\
R-squared Adj. & 0.122 & & \\
Log-Likelihood & -487.089 & -462.420 & -462.554 \\
Nobs & 753 & 753 & 753 \\
Pseudo R-squared & - & 0.102 & 0.102 \\
R-squared & 0.130 & - & - \\
\end{longtable}
}

\emph{Signif. codes:} \texttt{***} \(p<0.01\), \texttt{**} \(p<0.05\),
\texttt{*} \(p<0.1\)

Вычислите предельное значение для каждого регрессора в средней точке
(MEM) для logit модели. \textbf{Ответ округлите до 3-х десятичных
знаков.}

Ответ:

{\def\LTcaptype{none} % do not increment counter
\begin{longtable}[]{@{}
  >{\raggedright\arraybackslash}p{(\linewidth - 12\tabcolsep) * \real{0.1905}}
  >{\raggedleft\arraybackslash}p{(\linewidth - 12\tabcolsep) * \real{0.1071}}
  >{\raggedleft\arraybackslash}p{(\linewidth - 12\tabcolsep) * \real{0.1548}}
  >{\raggedleft\arraybackslash}p{(\linewidth - 12\tabcolsep) * \real{0.1310}}
  >{\raggedleft\arraybackslash}p{(\linewidth - 12\tabcolsep) * \real{0.1310}}
  >{\raggedleft\arraybackslash}p{(\linewidth - 12\tabcolsep) * \real{0.1429}}
  >{\raggedleft\arraybackslash}p{(\linewidth - 12\tabcolsep) * \real{0.1429}}@{}}
\toprule\noalign{}
\begin{minipage}[b]{\linewidth}\raggedright
\end{minipage} & \begin{minipage}[b]{\linewidth}\raggedleft
dy/dx
\end{minipage} & \begin{minipage}[b]{\linewidth}\raggedleft
Std. Err.
\end{minipage} & \begin{minipage}[b]{\linewidth}\raggedleft
z-value
\end{minipage} & \begin{minipage}[b]{\linewidth}\raggedleft
p-value
\end{minipage} & \begin{minipage}[b]{\linewidth}\raggedleft
CI Lower
\end{minipage} & \begin{minipage}[b]{\linewidth}\raggedleft
CI Upper
\end{minipage} \\
\midrule\noalign{}
\endhead
\bottomrule\noalign{}
\endlastfoot
WA & -0.015 & 0.003 & -5.006 & 0 & -0.021 & -0.009 \\
np.log(FAMINC) & 0.083 & 0.042 & 1.982 & 0.048 & 0.001 & 0.166 \\
WE & 0.044 & 0.01 & 4.45 & 0 & 0.025 & 0.063 \\
KL6 & -0.353 & 0.048 & -7.395 & 0 & -0.446 & -0.259 \\
K618 & -0.023 & 0.016 & -1.416 & 0.157 & -0.055 & 0.009 \\
CIT & -0.052 & 0.043 & -1.211 & 0.226 & -0.137 & 0.032 \\
UN & -0.004 & 0.006 & -0.675 & 0.5 & -0.017 & 0.008 \\
\end{longtable}
}

Вычислите предельное значение для каждого регрессора в средней точке
(MEM) для probit модели. \textbf{Ответ округлите до 3-х десятичных
знаков.}

Ответ:

{\def\LTcaptype{none} % do not increment counter
\begin{longtable}[]{@{}
  >{\raggedright\arraybackslash}p{(\linewidth - 12\tabcolsep) * \real{0.1905}}
  >{\raggedleft\arraybackslash}p{(\linewidth - 12\tabcolsep) * \real{0.1071}}
  >{\raggedleft\arraybackslash}p{(\linewidth - 12\tabcolsep) * \real{0.1548}}
  >{\raggedleft\arraybackslash}p{(\linewidth - 12\tabcolsep) * \real{0.1310}}
  >{\raggedleft\arraybackslash}p{(\linewidth - 12\tabcolsep) * \real{0.1310}}
  >{\raggedleft\arraybackslash}p{(\linewidth - 12\tabcolsep) * \real{0.1429}}
  >{\raggedleft\arraybackslash}p{(\linewidth - 12\tabcolsep) * \real{0.1429}}@{}}
\toprule\noalign{}
\begin{minipage}[b]{\linewidth}\raggedright
\end{minipage} & \begin{minipage}[b]{\linewidth}\raggedleft
dy/dx
\end{minipage} & \begin{minipage}[b]{\linewidth}\raggedleft
Std. Err.
\end{minipage} & \begin{minipage}[b]{\linewidth}\raggedleft
z-value
\end{minipage} & \begin{minipage}[b]{\linewidth}\raggedleft
p-value
\end{minipage} & \begin{minipage}[b]{\linewidth}\raggedleft
CI Lower
\end{minipage} & \begin{minipage}[b]{\linewidth}\raggedleft
CI Upper
\end{minipage} \\
\midrule\noalign{}
\endhead
\bottomrule\noalign{}
\endlastfoot
WA & -0.015 & 0.003 & -5.053 & 0 & -0.021 & -0.009 \\
np.log(FAMINC) & 0.081 & 0.041 & 1.973 & 0.049 & 0.001 & 0.161 \\
WE & 0.042 & 0.009 & 4.504 & 0 & 0.024 & 0.061 \\
KL6 & -0.34 & 0.044 & -7.738 & 0 & -0.427 & -0.254 \\
K618 & -0.022 & 0.016 & -1.407 & 0.16 & -0.053 & 0.009 \\
CIT & -0.049 & 0.042 & -1.17 & 0.242 & -0.132 & 0.033 \\
UN & -0.004 & 0.006 & -0.67 & 0.503 & -0.016 & 0.008 \\
\end{longtable}
}

Вычислите среднее по выборке предельное значение (AME) для каждого
регрессора для logit модели. \textbf{Ответ округлите до 3-х десятичных
знаков.}

Ответ:

{\def\LTcaptype{none} % do not increment counter
\begin{longtable}[]{@{}
  >{\raggedright\arraybackslash}p{(\linewidth - 12\tabcolsep) * \real{0.1905}}
  >{\raggedleft\arraybackslash}p{(\linewidth - 12\tabcolsep) * \real{0.1071}}
  >{\raggedleft\arraybackslash}p{(\linewidth - 12\tabcolsep) * \real{0.1548}}
  >{\raggedleft\arraybackslash}p{(\linewidth - 12\tabcolsep) * \real{0.1310}}
  >{\raggedleft\arraybackslash}p{(\linewidth - 12\tabcolsep) * \real{0.1310}}
  >{\raggedleft\arraybackslash}p{(\linewidth - 12\tabcolsep) * \real{0.1429}}
  >{\raggedleft\arraybackslash}p{(\linewidth - 12\tabcolsep) * \real{0.1429}}@{}}
\toprule\noalign{}
\begin{minipage}[b]{\linewidth}\raggedright
\end{minipage} & \begin{minipage}[b]{\linewidth}\raggedleft
dy/dx
\end{minipage} & \begin{minipage}[b]{\linewidth}\raggedleft
Std. Err.
\end{minipage} & \begin{minipage}[b]{\linewidth}\raggedleft
z-value
\end{minipage} & \begin{minipage}[b]{\linewidth}\raggedleft
p-value
\end{minipage} & \begin{minipage}[b]{\linewidth}\raggedleft
CI Lower
\end{minipage} & \begin{minipage}[b]{\linewidth}\raggedleft
CI Upper
\end{minipage} \\
\midrule\noalign{}
\endhead
\bottomrule\noalign{}
\endlastfoot
WA & -0.013 & 0.003 & -5.324 & 0 & -0.018 & -0.008 \\
np.log(FAMINC) & 0.073 & 0.036 & 2.001 & 0.045 & 0.001 & 0.144 \\
WE & 0.038 & 0.008 & 4.667 & 0 & 0.022 & 0.054 \\
KL6 & -0.307 & 0.036 & -8.642 & 0 & -0.377 & -0.238 \\
K618 & -0.02 & 0.014 & -1.422 & 0.155 & -0.048 & 0.008 \\
CIT & -0.046 & 0.037 & -1.215 & 0.224 & -0.119 & 0.028 \\
UN & -0.004 & 0.005 & -0.676 & 0.499 & -0.014 & 0.007 \\
\end{longtable}
}

Вычислите среднее по выборке предельное значение (AME) для каждого
регрессора для probit модели. \textbf{Ответ округлите до 3-х десятичных
знаков.}

Ответ:

{\def\LTcaptype{none} % do not increment counter
\begin{longtable}[]{@{}
  >{\raggedright\arraybackslash}p{(\linewidth - 12\tabcolsep) * \real{0.1905}}
  >{\raggedleft\arraybackslash}p{(\linewidth - 12\tabcolsep) * \real{0.1071}}
  >{\raggedleft\arraybackslash}p{(\linewidth - 12\tabcolsep) * \real{0.1548}}
  >{\raggedleft\arraybackslash}p{(\linewidth - 12\tabcolsep) * \real{0.1310}}
  >{\raggedleft\arraybackslash}p{(\linewidth - 12\tabcolsep) * \real{0.1310}}
  >{\raggedleft\arraybackslash}p{(\linewidth - 12\tabcolsep) * \real{0.1429}}
  >{\raggedleft\arraybackslash}p{(\linewidth - 12\tabcolsep) * \real{0.1429}}@{}}
\toprule\noalign{}
\begin{minipage}[b]{\linewidth}\raggedright
\end{minipage} & \begin{minipage}[b]{\linewidth}\raggedleft
dy/dx
\end{minipage} & \begin{minipage}[b]{\linewidth}\raggedleft
Std. Err.
\end{minipage} & \begin{minipage}[b]{\linewidth}\raggedleft
z-value
\end{minipage} & \begin{minipage}[b]{\linewidth}\raggedleft
p-value
\end{minipage} & \begin{minipage}[b]{\linewidth}\raggedleft
CI Lower
\end{minipage} & \begin{minipage}[b]{\linewidth}\raggedleft
CI Upper
\end{minipage} \\
\midrule\noalign{}
\endhead
\bottomrule\noalign{}
\endlastfoot
WA & -0.013 & 0.003 & -5.314 & 0 & -0.018 & -0.008 \\
np.log(FAMINC) & 0.072 & 0.036 & 1.987 & 0.047 & 0.001 & 0.143 \\
WE & 0.038 & 0.008 & 4.686 & 0 & 0.022 & 0.054 \\
KL6 & -0.304 & 0.035 & -8.818 & 0 & -0.372 & -0.237 \\
K618 & -0.02 & 0.014 & -1.412 & 0.158 & -0.048 & 0.008 \\
CIT & -0.044 & 0.038 & -1.173 & 0.241 & -0.118 & 0.03 \\
UN & -0.004 & 0.005 & -0.671 & 0.502 & -0.014 & 0.007 \\
\end{longtable}
}

\section{approve equation}\label{approve-equation}

Для датасета \texttt{loanapp} рассмотрим регрессию \textbf{approve на
I(appinc / 100), mortno, unem, dep, male} трёх спецификаций:

\begin{itemize}
\tightlist
\item
  LPM
\item
  logit
\item
  probit
\end{itemize}

Результаты подгонки:

{\def\LTcaptype{none} % do not increment counter
\begin{longtable}[]{@{}llll@{}}
\toprule\noalign{}
& LPM & logit & probit \\
\midrule\noalign{}
\endhead
\bottomrule\noalign{}
\endlastfoot
Dep. Variable & approve & approve & approve \\
I(appinc / 100) & -0.014* & -0.106 & -0.056 \\
& (0.009) & (0.067) & (0.038) \\
mortno & 0.079*** & 0.817*** & 0.422*** \\
& (0.016) & (0.170) & (0.086) \\
unem & -0.008** & -0.065** & -0.036** \\
& (0.003) & (0.029) & (0.016) \\
dep & -0.012* & -0.106* & -0.055* \\
& (0.007) & (0.061) & (0.033) \\
male & 0.019 & 0.173 & 0.093 \\
& (0.019) & (0.176) & (0.094) \\
Intercept & 0.886*** & 2.032*** & 1.198*** \\
& (0.022) & (0.193) & (0.105) \\
R-squared & 0.017 & & \\
R-squared Adj. & 0.014 & & \\
Log-Likelihood & -590.863 & -720.861 & -720.982 \\
Nobs & 1971 & 1971 & 1971 \\
Pseudo R-squared & - & 0.023 & 0.023 \\
R-squared & 0.017 & - & - \\
\end{longtable}
}

\emph{Signif. codes:} \texttt{***} \(p<0.01\), \texttt{**} \(p<0.05\),
\texttt{*} \(p<0.1\)

Вычислите предельное значение для каждого регрессора в средней точке
(MEM) для logit модели. \textbf{Ответ округлите до 3-х десятичных
знаков.}

Ответ:

{\def\LTcaptype{none} % do not increment counter
\begin{longtable}[]{@{}
  >{\raggedright\arraybackslash}p{(\linewidth - 12\tabcolsep) * \real{0.2000}}
  >{\raggedleft\arraybackslash}p{(\linewidth - 12\tabcolsep) * \real{0.1059}}
  >{\raggedleft\arraybackslash}p{(\linewidth - 12\tabcolsep) * \real{0.1529}}
  >{\raggedleft\arraybackslash}p{(\linewidth - 12\tabcolsep) * \real{0.1294}}
  >{\raggedleft\arraybackslash}p{(\linewidth - 12\tabcolsep) * \real{0.1294}}
  >{\raggedleft\arraybackslash}p{(\linewidth - 12\tabcolsep) * \real{0.1412}}
  >{\raggedleft\arraybackslash}p{(\linewidth - 12\tabcolsep) * \real{0.1412}}@{}}
\toprule\noalign{}
\begin{minipage}[b]{\linewidth}\raggedright
\end{minipage} & \begin{minipage}[b]{\linewidth}\raggedleft
dy/dx
\end{minipage} & \begin{minipage}[b]{\linewidth}\raggedleft
Std. Err.
\end{minipage} & \begin{minipage}[b]{\linewidth}\raggedleft
z-value
\end{minipage} & \begin{minipage}[b]{\linewidth}\raggedleft
p-value
\end{minipage} & \begin{minipage}[b]{\linewidth}\raggedleft
CI Lower
\end{minipage} & \begin{minipage}[b]{\linewidth}\raggedleft
CI Upper
\end{minipage} \\
\midrule\noalign{}
\endhead
\bottomrule\noalign{}
\endlastfoot
I(appinc / 100) & -0.011 & 0.007 & -1.583 & 0.114 & -0.025 & 0.003 \\
mortno & 0.084 & 0.017 & 5 & 0 & 0.051 & 0.118 \\
unem & -0.007 & 0.003 & -2.26 & 0.024 & -0.013 & -0.001 \\
dep & -0.011 & 0.006 & -1.741 & 0.082 & -0.023 & 0.001 \\
male & 0.018 & 0.018 & 0.985 & 0.325 & -0.018 & 0.054 \\
\end{longtable}
}

Вычислите предельное значение для каждого регрессора в средней точке
(MEM) для probit модели. \textbf{Ответ округлите до 3-х десятичных
знаков.}

Ответ:

{\def\LTcaptype{none} % do not increment counter
\begin{longtable}[]{@{}
  >{\raggedright\arraybackslash}p{(\linewidth - 12\tabcolsep) * \real{0.2000}}
  >{\raggedleft\arraybackslash}p{(\linewidth - 12\tabcolsep) * \real{0.1059}}
  >{\raggedleft\arraybackslash}p{(\linewidth - 12\tabcolsep) * \real{0.1529}}
  >{\raggedleft\arraybackslash}p{(\linewidth - 12\tabcolsep) * \real{0.1294}}
  >{\raggedleft\arraybackslash}p{(\linewidth - 12\tabcolsep) * \real{0.1294}}
  >{\raggedleft\arraybackslash}p{(\linewidth - 12\tabcolsep) * \real{0.1412}}
  >{\raggedleft\arraybackslash}p{(\linewidth - 12\tabcolsep) * \real{0.1412}}@{}}
\toprule\noalign{}
\begin{minipage}[b]{\linewidth}\raggedright
\end{minipage} & \begin{minipage}[b]{\linewidth}\raggedleft
dy/dx
\end{minipage} & \begin{minipage}[b]{\linewidth}\raggedleft
Std. Err.
\end{minipage} & \begin{minipage}[b]{\linewidth}\raggedleft
z-value
\end{minipage} & \begin{minipage}[b]{\linewidth}\raggedleft
p-value
\end{minipage} & \begin{minipage}[b]{\linewidth}\raggedleft
CI Lower
\end{minipage} & \begin{minipage}[b]{\linewidth}\raggedleft
CI Upper
\end{minipage} \\
\midrule\noalign{}
\endhead
\bottomrule\noalign{}
\endlastfoot
I(appinc / 100) & -0.011 & 0.008 & -1.487 & 0.137 & -0.026 & 0.004 \\
mortno & 0.084 & 0.017 & 5.014 & 0 & 0.051 & 0.116 \\
unem & -0.007 & 0.003 & -2.281 & 0.023 & -0.013 & -0.001 \\
dep & -0.011 & 0.007 & -1.65 & 0.099 & -0.024 & 0.002 \\
male & 0.018 & 0.019 & 0.987 & 0.324 & -0.018 & 0.055 \\
\end{longtable}
}

Вычислите среднее по выборке предельное значение (AME) для каждого
регрессора для logit модели. \textbf{Ответ округлите до 3-х десятичных
знаков.}

Ответ:

{\def\LTcaptype{none} % do not increment counter
\begin{longtable}[]{@{}
  >{\raggedright\arraybackslash}p{(\linewidth - 12\tabcolsep) * \real{0.2000}}
  >{\raggedleft\arraybackslash}p{(\linewidth - 12\tabcolsep) * \real{0.1059}}
  >{\raggedleft\arraybackslash}p{(\linewidth - 12\tabcolsep) * \real{0.1529}}
  >{\raggedleft\arraybackslash}p{(\linewidth - 12\tabcolsep) * \real{0.1294}}
  >{\raggedleft\arraybackslash}p{(\linewidth - 12\tabcolsep) * \real{0.1294}}
  >{\raggedleft\arraybackslash}p{(\linewidth - 12\tabcolsep) * \real{0.1412}}
  >{\raggedleft\arraybackslash}p{(\linewidth - 12\tabcolsep) * \real{0.1412}}@{}}
\toprule\noalign{}
\begin{minipage}[b]{\linewidth}\raggedright
\end{minipage} & \begin{minipage}[b]{\linewidth}\raggedleft
dy/dx
\end{minipage} & \begin{minipage}[b]{\linewidth}\raggedleft
Std. Err.
\end{minipage} & \begin{minipage}[b]{\linewidth}\raggedleft
z-value
\end{minipage} & \begin{minipage}[b]{\linewidth}\raggedleft
p-value
\end{minipage} & \begin{minipage}[b]{\linewidth}\raggedleft
CI Lower
\end{minipage} & \begin{minipage}[b]{\linewidth}\raggedleft
CI Upper
\end{minipage} \\
\midrule\noalign{}
\endhead
\bottomrule\noalign{}
\endlastfoot
I(appinc / 100) & -0.011 & 0.007 & -1.581 & 0.114 & -0.025 & 0.003 \\
mortno & 0.087 & 0.018 & 4.764 & 0 & 0.051 & 0.123 \\
unem & -0.007 & 0.003 & -2.251 & 0.024 & -0.013 & -0.001 \\
dep & -0.011 & 0.007 & -1.736 & 0.083 & -0.024 & 0.001 \\
male & 0.018 & 0.019 & 0.984 & 0.325 & -0.018 & 0.055 \\
\end{longtable}
}

Вычислите среднее по выборке предельное значение (AME) для каждого
регрессора для probit модели. \textbf{Ответ округлите до 3-х десятичных
знаков.}

Ответ:

{\def\LTcaptype{none} % do not increment counter
\begin{longtable}[]{@{}
  >{\raggedright\arraybackslash}p{(\linewidth - 12\tabcolsep) * \real{0.2000}}
  >{\raggedleft\arraybackslash}p{(\linewidth - 12\tabcolsep) * \real{0.1059}}
  >{\raggedleft\arraybackslash}p{(\linewidth - 12\tabcolsep) * \real{0.1529}}
  >{\raggedleft\arraybackslash}p{(\linewidth - 12\tabcolsep) * \real{0.1294}}
  >{\raggedleft\arraybackslash}p{(\linewidth - 12\tabcolsep) * \real{0.1294}}
  >{\raggedleft\arraybackslash}p{(\linewidth - 12\tabcolsep) * \real{0.1412}}
  >{\raggedleft\arraybackslash}p{(\linewidth - 12\tabcolsep) * \real{0.1412}}@{}}
\toprule\noalign{}
\begin{minipage}[b]{\linewidth}\raggedright
\end{minipage} & \begin{minipage}[b]{\linewidth}\raggedleft
dy/dx
\end{minipage} & \begin{minipage}[b]{\linewidth}\raggedleft
Std. Err.
\end{minipage} & \begin{minipage}[b]{\linewidth}\raggedleft
z-value
\end{minipage} & \begin{minipage}[b]{\linewidth}\raggedleft
p-value
\end{minipage} & \begin{minipage}[b]{\linewidth}\raggedleft
CI Lower
\end{minipage} & \begin{minipage}[b]{\linewidth}\raggedleft
CI Upper
\end{minipage} \\
\midrule\noalign{}
\endhead
\bottomrule\noalign{}
\endlastfoot
I(appinc / 100) & -0.011 & 0.008 & -1.488 & 0.137 & -0.026 & 0.004 \\
mortno & 0.084 & 0.017 & 4.931 & 0 & 0.051 & 0.118 \\
unem & -0.007 & 0.003 & -2.28 & 0.023 & -0.013 & -0.001 \\
dep & -0.011 & 0.007 & -1.649 & 0.099 & -0.024 & 0.002 \\
male & 0.019 & 0.019 & 0.987 & 0.324 & -0.018 & 0.056 \\
\end{longtable}
}

\section{swiss labour force equation}\label{swiss-labour-force-equation}

Для датасета \texttt{SwissLabour} рассмотрим регрессию
\textbf{participation на income, age, I(age ** ), youngkids, oldkids}
трёх спецификаций:

\begin{itemize}
\tightlist
\item
  LPM
\item
  logit
\item
  probit
\end{itemize}

Результаты подгонки:

{\def\LTcaptype{none} % do not increment counter
\begin{longtable}[]{@{}
  >{\raggedright\arraybackslash}p{(\linewidth - 6\tabcolsep) * \real{0.2400}}
  >{\raggedright\arraybackslash}p{(\linewidth - 6\tabcolsep) * \real{0.2533}}
  >{\raggedright\arraybackslash}p{(\linewidth - 6\tabcolsep) * \real{0.2533}}
  >{\raggedright\arraybackslash}p{(\linewidth - 6\tabcolsep) * \real{0.2533}}@{}}
\toprule\noalign{}
\begin{minipage}[b]{\linewidth}\raggedright
\end{minipage} & \begin{minipage}[b]{\linewidth}\raggedright
LPM
\end{minipage} & \begin{minipage}[b]{\linewidth}\raggedright
logit
\end{minipage} & \begin{minipage}[b]{\linewidth}\raggedright
probit
\end{minipage} \\
\midrule\noalign{}
\endhead
\bottomrule\noalign{}
\endlastfoot
Dep. Variable & participation\_num & participation\_num &
participation\_num \\
income & -0.258*** & -1.337*** & -0.805*** \\
& (0.039) & (0.214) & (0.125) \\
age & 0.795*** & 3.907*** & 2.361*** \\
& (0.131) & (0.670) & (0.397) \\
I(age ** 2) & -0.112*** & -0.550*** & -0.332*** \\
& (0.016) & (0.083) & (0.049) \\
youngkids & -0.228*** & -1.069*** & -0.640*** \\
& (0.032) & (0.166) & (0.096) \\
oldkids & -0.053*** & -0.252*** & -0.152*** \\
& (0.017) & (0.083) & (0.050) \\
Intercept & 2.082*** & 8.431*** & 5.055*** \\
& (0.446) & (2.330) & (1.388) \\
R-squared & 0.155 & & \\
R-squared Adj. & 0.150 & & \\
Log-Likelihood & -556.793 & -527.064 & -527.405 \\
Nobs & 872 & 872 & 872 \\
Pseudo R-squared & - & 0.124 & 0.123 \\
R-squared & 0.155 & - & - \\
\end{longtable}
}

\emph{Signif. codes:} \texttt{***} \(p<0.01\), \texttt{**} \(p<0.05\),
\texttt{*} \(p<0.1\)

Вычислите предельное значение для каждого регрессора в средней точке
(MEM) для logit модели. \textbf{Ответ округлите до 3-х десятичных
знаков.}

Ответ:

{\def\LTcaptype{none} % do not increment counter
\begin{longtable}[]{@{}
  >{\raggedright\arraybackslash}p{(\linewidth - 12\tabcolsep) * \real{0.1605}}
  >{\raggedleft\arraybackslash}p{(\linewidth - 12\tabcolsep) * \real{0.1111}}
  >{\raggedleft\arraybackslash}p{(\linewidth - 12\tabcolsep) * \real{0.1605}}
  >{\raggedleft\arraybackslash}p{(\linewidth - 12\tabcolsep) * \real{0.1358}}
  >{\raggedleft\arraybackslash}p{(\linewidth - 12\tabcolsep) * \real{0.1358}}
  >{\raggedleft\arraybackslash}p{(\linewidth - 12\tabcolsep) * \real{0.1481}}
  >{\raggedleft\arraybackslash}p{(\linewidth - 12\tabcolsep) * \real{0.1481}}@{}}
\toprule\noalign{}
\begin{minipage}[b]{\linewidth}\raggedright
\end{minipage} & \begin{minipage}[b]{\linewidth}\raggedleft
dy/dx
\end{minipage} & \begin{minipage}[b]{\linewidth}\raggedleft
Std. Err.
\end{minipage} & \begin{minipage}[b]{\linewidth}\raggedleft
z-value
\end{minipage} & \begin{minipage}[b]{\linewidth}\raggedleft
p-value
\end{minipage} & \begin{minipage}[b]{\linewidth}\raggedleft
CI Lower
\end{minipage} & \begin{minipage}[b]{\linewidth}\raggedleft
CI Upper
\end{minipage} \\
\midrule\noalign{}
\endhead
\bottomrule\noalign{}
\endlastfoot
income & -0.33 & 0.053 & -6.249 & 0 & -0.434 & -0.227 \\
age & 0.965 & 0.165 & 5.849 & 0 & 0.642 & 1.289 \\
I(age ** 2) & -0.136 & 0.02 & -6.633 & 0 & -0.176 & -0.096 \\
youngkids & -0.264 & 0.041 & -6.458 & 0 & -0.344 & -0.184 \\
oldkids & -0.062 & 0.02 & -3.052 & 0.002 & -0.102 & -0.022 \\
\end{longtable}
}

Вычислите предельное значение для каждого регрессора в средней точке
(MEM) для probit модели. \textbf{Ответ округлите до 3-х десятичных
знаков.}

Ответ:

{\def\LTcaptype{none} % do not increment counter
\begin{longtable}[]{@{}
  >{\raggedright\arraybackslash}p{(\linewidth - 12\tabcolsep) * \real{0.1605}}
  >{\raggedleft\arraybackslash}p{(\linewidth - 12\tabcolsep) * \real{0.1111}}
  >{\raggedleft\arraybackslash}p{(\linewidth - 12\tabcolsep) * \real{0.1605}}
  >{\raggedleft\arraybackslash}p{(\linewidth - 12\tabcolsep) * \real{0.1358}}
  >{\raggedleft\arraybackslash}p{(\linewidth - 12\tabcolsep) * \real{0.1358}}
  >{\raggedleft\arraybackslash}p{(\linewidth - 12\tabcolsep) * \real{0.1481}}
  >{\raggedleft\arraybackslash}p{(\linewidth - 12\tabcolsep) * \real{0.1481}}@{}}
\toprule\noalign{}
\begin{minipage}[b]{\linewidth}\raggedright
\end{minipage} & \begin{minipage}[b]{\linewidth}\raggedleft
dy/dx
\end{minipage} & \begin{minipage}[b]{\linewidth}\raggedleft
Std. Err.
\end{minipage} & \begin{minipage}[b]{\linewidth}\raggedleft
z-value
\end{minipage} & \begin{minipage}[b]{\linewidth}\raggedleft
p-value
\end{minipage} & \begin{minipage}[b]{\linewidth}\raggedleft
CI Lower
\end{minipage} & \begin{minipage}[b]{\linewidth}\raggedleft
CI Upper
\end{minipage} \\
\midrule\noalign{}
\endhead
\bottomrule\noalign{}
\endlastfoot
income & -0.318 & 0.05 & -6.428 & 0 & -0.415 & -0.221 \\
age & 0.934 & 0.157 & 5.95 & 0 & 0.627 & 1.242 \\
I(age ** 2) & -0.131 & 0.019 & -6.79 & 0 & -0.169 & -0.093 \\
youngkids & -0.253 & 0.038 & -6.69 & 0 & -0.327 & -0.179 \\
oldkids & -0.06 & 0.02 & -3.055 & 0.002 & -0.099 & -0.022 \\
\end{longtable}
}

Вычислите среднее по выборке предельное значение (AME) для каждого
регрессора для logit модели. \textbf{Ответ округлите до 3-х десятичных
знаков.}

Ответ:

{\def\LTcaptype{none} % do not increment counter
\begin{longtable}[]{@{}
  >{\raggedright\arraybackslash}p{(\linewidth - 12\tabcolsep) * \real{0.1605}}
  >{\raggedleft\arraybackslash}p{(\linewidth - 12\tabcolsep) * \real{0.1111}}
  >{\raggedleft\arraybackslash}p{(\linewidth - 12\tabcolsep) * \real{0.1605}}
  >{\raggedleft\arraybackslash}p{(\linewidth - 12\tabcolsep) * \real{0.1358}}
  >{\raggedleft\arraybackslash}p{(\linewidth - 12\tabcolsep) * \real{0.1358}}
  >{\raggedleft\arraybackslash}p{(\linewidth - 12\tabcolsep) * \real{0.1481}}
  >{\raggedleft\arraybackslash}p{(\linewidth - 12\tabcolsep) * \real{0.1481}}@{}}
\toprule\noalign{}
\begin{minipage}[b]{\linewidth}\raggedright
\end{minipage} & \begin{minipage}[b]{\linewidth}\raggedleft
dy/dx
\end{minipage} & \begin{minipage}[b]{\linewidth}\raggedleft
Std. Err.
\end{minipage} & \begin{minipage}[b]{\linewidth}\raggedleft
z-value
\end{minipage} & \begin{minipage}[b]{\linewidth}\raggedleft
p-value
\end{minipage} & \begin{minipage}[b]{\linewidth}\raggedleft
CI Lower
\end{minipage} & \begin{minipage}[b]{\linewidth}\raggedleft
CI Upper
\end{minipage} \\
\midrule\noalign{}
\endhead
\bottomrule\noalign{}
\endlastfoot
income & -0.279 & 0.041 & -6.788 & 0 & -0.36 & -0.199 \\
age & 0.817 & 0.13 & 6.288 & 0 & 0.562 & 1.071 \\
I(age ** 2) & -0.115 & 0.016 & -7.295 & 0 & -0.146 & -0.084 \\
youngkids & -0.224 & 0.032 & -7.082 & 0 & -0.285 & -0.162 \\
oldkids & -0.053 & 0.017 & -3.112 & 0.002 & -0.086 & -0.019 \\
\end{longtable}
}

Вычислите среднее по выборке предельное значение (AME) для каждого
регрессора для probit модели. \textbf{Ответ округлите до 3-х десятичных
знаков.}

Ответ:

{\def\LTcaptype{none} % do not increment counter
\begin{longtable}[]{@{}
  >{\raggedright\arraybackslash}p{(\linewidth - 12\tabcolsep) * \real{0.1605}}
  >{\raggedleft\arraybackslash}p{(\linewidth - 12\tabcolsep) * \real{0.1111}}
  >{\raggedleft\arraybackslash}p{(\linewidth - 12\tabcolsep) * \real{0.1605}}
  >{\raggedleft\arraybackslash}p{(\linewidth - 12\tabcolsep) * \real{0.1358}}
  >{\raggedleft\arraybackslash}p{(\linewidth - 12\tabcolsep) * \real{0.1358}}
  >{\raggedleft\arraybackslash}p{(\linewidth - 12\tabcolsep) * \real{0.1481}}
  >{\raggedleft\arraybackslash}p{(\linewidth - 12\tabcolsep) * \real{0.1481}}@{}}
\toprule\noalign{}
\begin{minipage}[b]{\linewidth}\raggedright
\end{minipage} & \begin{minipage}[b]{\linewidth}\raggedleft
dy/dx
\end{minipage} & \begin{minipage}[b]{\linewidth}\raggedleft
Std. Err.
\end{minipage} & \begin{minipage}[b]{\linewidth}\raggedleft
z-value
\end{minipage} & \begin{minipage}[b]{\linewidth}\raggedleft
p-value
\end{minipage} & \begin{minipage}[b]{\linewidth}\raggedleft
CI Lower
\end{minipage} & \begin{minipage}[b]{\linewidth}\raggedleft
CI Upper
\end{minipage} \\
\midrule\noalign{}
\endhead
\bottomrule\noalign{}
\endlastfoot
income & -0.278 & 0.04 & -6.898 & 0 & -0.357 & -0.199 \\
age & 0.815 & 0.129 & 6.323 & 0 & 0.562 & 1.067 \\
I(age ** 2) & -0.115 & 0.016 & -7.359 & 0 & -0.145 & -0.084 \\
youngkids & -0.221 & 0.031 & -7.232 & 0 & -0.281 & -0.161 \\
oldkids & -0.053 & 0.017 & -3.104 & 0.002 & -0.086 & -0.019 \\
\end{longtable}
}




\end{document}
